\documentclass[11pt,a4paper]{article}

\usepackage[margin=1.05in]{geometry}
\usepackage{amsmath,amssymb,amsthm}
\usepackage{mathtools}
\usepackage{stmaryrd}
\usepackage{microtype}
\setlength{\emergencystretch}{2.5em}
\usepackage{enumitem}
\usepackage{booktabs}
\usepackage{array}
\usepackage{xcolor}
\usepackage[colorlinks=true,linkcolor=blue!45!black,urlcolor=blue!60!black,citecolor=blue!45!black]{hyperref}

\newtheorem{theorem}{Theorem}[section]
\newtheorem{proposition}[theorem]{Proposition}
\newtheorem{lemma}[theorem]{Lemma}
\newtheorem{corollary}[theorem]{Corollary}
\newtheorem{conjecture}[theorem]{Conjecture}
\theoremstyle{definition}
\newtheorem{definition}[theorem]{Definition}
\newtheorem{example}[theorem]{Example}
\theoremstyle{remark}
\newtheorem{remark}[theorem]{Remark}
\newtheorem{warning}[theorem]{Warning}
\newtheorem{question}[theorem]{Question}

\newcommand{\FF}{\mathbb{F}}
\newcommand{\NN}{\mathbb{N}}
\newcommand{\ZZ}{\mathbb{Z}}
\newcommand{\RR}{\mathbb{R}}
\newcommand{\cat}[1]{\mathbf{#1}}
\newcommand{\Set}{\cat{Set}}
\newcommand{\FinSet}{\cat{FinSet}}
\newcommand{\Poly}{\cat{Poly}}
\newcommand{\Cont}{\cat{Cont}}
\newcommand{\vect}{\cat{vect}}
\newcommand{\iso}{\mathrm{iso}}
\newcommand{\con}{\mathrel{\vartriangleright}}
\newcommand{\norm}[1]{\lvert #1 \rvert}
\DeclareMathOperator{\Aut}{Aut}
\DeclareMathOperator{\GL}{GL}
\DeclareMathOperator{\Bag}{Bag}
\DeclareMathOperator{\List}{List}

\definecolor{okgreen}{rgb}{0.0,0.45,0.15}
\definecolor{nored}{rgb}{0.65,0.0,0.0}
\newcommand{\OK}{\textcolor{okgreen}{\textbf{verified}}}
\newcommand{\NO}{\textcolor{nored}{\textbf{fails}}}

\title{\bfseries List and Bag, $\pi$ and $\sqrt{\pi}$\\[0.4em]
\large The missing ordinary row of the decategorification table,
and a Gaussian functional that counts $S_n$-quotients}
\author{}
\date{}

\begin{document}
\maketitle

\begin{abstract}
The ordinary generating series $O(x) = \sum x^n$ and the exponential one
$E(x) = \sum x^n/n!$ satisfy
\[
  \int_{\RR} O(-x^2)\,dx = \pi,
  \qquad
  \int_{\RR} E(-x^2)\,dx = \sqrt{\pi}.
\]
Both are $\Gamma(\tfrac12)$ to a power. We read this containerwise. $O$ and $E$ are
not two generating functions of one species but the exponential generating functions
of two species, $L$ (linear orders) and $E$ (sets), which as containers are
$\List$ and $\Bag$: the same shapes and positions, with and without the
$S_n$-quotient. The functional $F \mapsto \int_{\RR} F(-x^2)\,dx$ sees the quotient,
and Theorem~\ref{thm:main} says exactly what it charges for it: dividing the
coefficients by $n!$ once costs one factor of $\Gamma(\tfrac12)$. The exponent $2$
in the $\pi$ case and the $-s$ in the general case have different origins --- the
first is the Gamma reflection formula at the self-dual Mellin exponent $\tfrac12$,
which is what the squaring in $-x^2$ is for, and the second is the $n!$-divisions.
Two obstructions are recorded honestly: $\Poly$ has no additive inverses, so $-x^2$
is not an object and the computation lives in the Grothendieck ring of virtual
species; and $\int_\RR$ is not a container operation but a second decategorification,
whose nearest precedent is Spivak's entropy of a polynomial functor.
\end{abstract}

\section{The observation}

\begin{proposition}\label{prop:obs}
$\displaystyle\int_{\RR}\frac{dx}{1+x^2} = \pi$ and
$\displaystyle\int_{\RR} e^{-x^2}\,dx = \sqrt{\pi}$.
\end{proposition}

Both are elementary; $O(-x^2) = 1/(1+x^2)$ because $O(x) = 1/(1-x)$, and
$E(-x^2) = e^{-x^2}$ because $E(x) = e^x$. Checked numerically to $30$ digits:
\OK. The integral must be taken over all of $\RR$; over $[0,\infty)$ one gets
$\pi/2$ and $\sqrt{\pi}/2$ and the pattern is obscured.

\begin{remark}[the termwise reading fails, and that is the content]\label{rem:termwise}
One cannot obtain either number by integrating the series. $\int_\RR x^{2n}\,dx = \infty$
for every $n$, so no rearrangement is available. Worse, for $O$ the series
$\sum (-x^2)^n$ converges only for $\norm{x} < 1$ while $1/(1+x^2)$ is integrable on
all of $\RR$: the functional is applied to the analytic continuation, not to the series.
The generating function is doing work here beyond bookkeeping, which is why the
identity is a statement about the \emph{middle rung} of the decategorification
diagram and not about coefficients.
\end{remark}

\section{The missing row}

Day convolution along a monoidal index category $(\mathcal{A}, \otimes)$ is the
unique monoidal structure making the Yoneda embedding strong monoidal, and
$\mathcal{A}$ is an input. Varying it varies the decategorified product. The table
of \texttt{q-dirichlet.tex} \S2 reads:

\begin{center}
\small
\begin{tabular}{@{}lll@{}}
\toprule
$\mathcal{A}$ & Day convolution & decategorification \\
\midrule
$(\FinSet^{\iso}, \sqcup)$ & species product & $\sum a_n x^n/n!$ (Cauchy/EGF) \\
$(\NN_{>0}, \times)$ & Dirichlet convolution & $\sum a_n n^{-s}$ \\
$(\FinSet^{\iso}_{\neq\emptyset}, \times)$ & arithmetic product & $\sum a_n n^{-s}/n!$ \\
\bottomrule
\end{tabular}
\end{center}

\noindent
The row the present question is about is not there. Let $\NN^{\delta}$ be the
\emph{discrete} monoid $(\NN, +)$, that is $\FinSet^{\iso}$ with every
non-identity automorphism deleted.

\begin{center}
\small
\begin{tabular}{@{}lll@{}}
\toprule
$\mathcal{A}$ & Day convolution & decategorification \\
\midrule
$(\NN^{\delta}, +)$ & Cauchy product of $\Set$-sequences & $\sum a_n x^n$ (OGF) \\
\bottomrule
\end{tabular}
\end{center}

\begin{proposition}[\emph{ours}]\label{prop:row}
The two rows $(\NN^{\delta},+)$ and $(\FinSet^{\iso},\sqcup)$ differ only in whether
the index category remembers $\Aut([n]) = S_n$, and the two decategorifications
differ only by the divisor $n! = \norm{S_n}$. Hence $O$ and $E$ of
Proposition~\ref{prop:obs} are one sequence of shapes, $a_n = 1$, decategorified
along two index categories related by the inclusion
$\NN^{\delta} \hookrightarrow \FinSet^{\iso}$.
\end{proposition}

This is the correct statement but it is not the most useful one, because it makes the
difference live in the \emph{functor}
$\mathcal{A} \mapsto {}$ (series ring) rather than in the objects. The next section
moves the difference into the containers, where it can be seen.

\section{$O$ and $E$ are $\List$ and $\Bag$}

Recall the two species of Joyal \cite{Joyal81}, in the notation of
Bergeron--Labelle--Leroux \cite{BLL}:
\[
  L[n] = \{\text{linear orders on } [n]\},\quad \norm{L[n]} = n!,
  \qquad
  E[n] = \{\ast\},\quad \norm{E[n]} = 1 .
\]
Their exponential generating functions are
\[
  \sum_{n\ge0} \norm{L[n]}\frac{x^n}{n!} = \sum_{n\ge 0} x^n = \frac{1}{1-x} = O(x),
  \qquad
  \sum_{n\ge0} \norm{E[n]}\frac{x^n}{n!} = e^x = E(x).
\]

\begin{proposition}[\emph{ours}, and immediate]\label{prop:listbag}
$O$ and $E$ are both exponential generating functions. $O$ is the EGF of the species
of linear orders; $E$ is the EGF of the species of sets. The ``ordinary versus
exponential'' distinction of \S2 is the same distinction as ``$L$ versus $E$'', seen
from the other side.
\end{proposition}

Containerwise, with $(S \con P)$ denoting shapes $S$ and positions $P$:
\[
  \List = \bigl(\NN \con n \mapsto [n]\bigr), \qquad
  \llbracket \List \rrbracket X = \sum_{n\ge 0} X^{n},
\]
the free monoid, and $\Bag$ is the same container with the $S_n$-action on positions
divided out,
\[
  \llbracket \Bag \rrbracket X = \sum_{n \ge 0} X^{n}/S_n ,
\]
a \emph{quotient container} in the sense of Kock \cite{Kock}: a polynomial functor
over groupoids rather than over sets. So:

\begin{center}
\small
\begin{tabular}{@{}llll@{}}
\toprule
container & species & EGF & $\int_\RR \cdot(-x^2)$ \\
\midrule
$\List$ & $L$, linear orders & $1/(1-x)$ & $\pi = \Gamma(\tfrac12)^2$ \\
$\Bag$ & $E$, sets & $e^x$ & $\sqrt{\pi} = \Gamma(\tfrac12)^1$ \\
\bottomrule
\end{tabular}
\end{center}

\noindent
This is the row of \texttt{q-container.tex} \S``Related work'' that names Kock's
species/quotient-container bridge as ``the structural model to imitate''. The
present note is a small instance of that bridge in which the imitation can be
carried out completely, because both sides are one-dimensional.

\section{Borel is the de-quotienting map}

The reason the two containers are comparable at all is that $\List$'s series is the
Laplace transform of $\Bag$'s. For $\mathrm{Re}\,x < 1$,
\begin{equation}\label{eq:borel}
  O(x) \;=\; \frac{1}{1-x} \;=\; \int_0^\infty e^{-t} e^{xt}\,dt \;=\; \int_0^\infty e^{-t}E(xt)\,dt,
\end{equation}
and the mechanism is coefficientwise: $\int_0^\infty e^{-t}t^n\,dt = n! = \norm{L[n]}$.
The Gamma integral is precisely what restores the $S_n$ that $\Bag$ divided out.

\begin{proposition}[\emph{ours}]\label{prop:two-gammas}
Substituting $x \mapsto -x^2$ in \eqref{eq:borel} and exchanging the order of
integration gives
\[
  \int_\RR O(-x^2)\,dx
  = \int_0^\infty e^{-t}\Bigl(\int_\RR e^{-tx^2}\,dx\Bigr) dt
  = \int_0^\infty e^{-t}\sqrt{\pi/t}\,\,dt
  = \sqrt{\pi}\,\Gamma(\tfrac12)
  = \Gamma(\tfrac12)^2 .
\]
\end{proposition}

The two factors of $\Gamma(\tfrac12)$ have visibly different sources: the inner
Gaussian supplies one, and the $t^{-1/2}$ left behind by the Gaussian scaling
$x \mapsto x/\sqrt{t}$ is eaten by the Borel integral to supply the other. So the
$\pi$ of the $\List$ row is the $\sqrt{\pi}$ of the $\Bag$ row multiplied by the cost
of de-quotienting.

\section{The theorem}

Proposition~\ref{prop:two-gammas} is one step of a family. Let
\[
  F_s(z) \;=\; \sum_{n \ge 0} \frac{z^n}{\Gamma(n+1)^s}
  \;=\; \sum_{n\ge0}\frac{z^n}{(n!)^s},
  \qquad s \in \ZZ_{\ge 0},
\]
so that $F_0 = O$, $F_1 = E$, and $F_2(z) = {}_0F_1(;1;z)$ with
$F_2(-x^2) = J_0(2x)$. Dividing by $n!$ once more is one more $S_n$-quotient taken.

\begin{theorem}[\emph{ours}]\label{thm:main}
For $m \ge 2$ and $s \ge 0$,
\[
  \int_0^\infty F_s(-x^m)\,dx
  \;=\; \frac{\pi}{m\sin(\pi/m)}\cdot\Gamma\!\bigl(1 - \tfrac1m\bigr)^{-s},
\]
in the sense of the Mellin regularisation below. In particular at $m = 2$,
\[
  \boxed{\;\int_{\RR} F_s(-x^2)\,dx \;=\; \Gamma(\tfrac12)\,\Gamma(\tfrac12)\cdot\Gamma(\tfrac12)^{-s}
  \;=\; \Gamma(\tfrac12)^{\,2-s} \;=\; \pi^{(2-s)/2}.\;}
\]
\end{theorem}

\begin{proof}
Ramanujan's master theorem in the form
$\int_0^\infty x^{\sigma-1}\sum_{k\ge0}(-1)^k\varphi(k)x^k\,dx = \Gamma(\sigma)\Gamma(1-\sigma)\varphi(-\sigma)$,
valid for $\varphi$ analytic in a right half-plane of suitable growth and
$0 < \mathrm{Re}\,\sigma < 1$. Substitute $u = x^m$, so $dx = \tfrac1m u^{1/m-1}du$:
\[
  \int_0^\infty F_s(-x^m)\,dx = \frac1m\int_0^\infty u^{\sigma-1}F_s(-u)\,du,
  \qquad \sigma = \tfrac1m \in (0,1).
\]
Here $\varphi(k) = \Gamma(k+1)^{-s}$, which is entire since $1/\Gamma$ is entire, so
$\varphi(-\sigma) = \Gamma(1-\tfrac1m)^{-s}$. The reflection formula gives
$\Gamma(\sigma)\Gamma(1-\sigma) = \pi/\sin(\pi\sigma)$, and the stated formula follows.
At $m = 2$ we have $\sigma = \tfrac12$, $\pi/\sin(\pi/2) = \pi = \Gamma(\tfrac12)^2$,
and $\varphi(-\tfrac12) = \Gamma(\tfrac12)^{-s}$; doubling for the two-sided integral
cancels the $1/m = \tfrac12$.
\end{proof}

\begin{remark}[what is special about $m = 2$, i.e.\ about the \emph{pair}]\label{rem:m2}
The two $\Gamma$-factors in Theorem~\ref{thm:main} are $\Gamma(\tfrac1m)\Gamma(1-\tfrac1m)$
from the reflection formula and $\Gamma(1-\tfrac1m)^{-s}$ from the $n!$-divisions.
At $m = 2$, and only there, $\tfrac1m = 1 - \tfrac1m$: the self-dual point of
$\sigma \mapsto 1-\sigma$. That is the unique exponent at which the whole answer is a
power of a single number rather than a product of two unrelated ones --- and hence
the unique one at which the pattern ``$\pi$, then $\sqrt{\pi}$'' is visible at all.
Containerwise, $x^2$ is $y \otimes y$, the \emph{pair}: substituting the pair container
is what places the Mellin exponent at the self-dual point.
\end{remark}

\begin{remark}[verification, and where convergence stops]\label{rem:verify}
Numerically (\texttt{mpmath}, 25--30 digits): $s = 0$ gives $\pi$ \OK;
$s = 1$ gives $\sqrt{\pi}$ \OK; $s = 2$ gives $1 = \Gamma(\tfrac12)^0$ \OK, as
$\int_\RR J_0(2x)\,dx = 1$, Abel-convergent and confirmed by damping
$e^{-\varepsilon x}$ with $\varepsilon = 2^{-1},\dots,2^{-5}$
($0.9701, 0.9923, 0.9981, 0.9995, 0.9999$). For $m = 3$, $s = 1$ the formula reduces
to $\Gamma(4/3) = \Gamma(1/3)/3$, confirmed symbolically by the reflection formula \OK.

For $s \ge 3$ the integral is \emph{not} absolutely convergent and need not be
conditionally convergent either: ${}_0F_{s-1}(;1,\dots,1;-x^2)$ oscillates in
$x^{2/s}$, so it is not periodic and ordinary oscillatory quadrature does not apply.
The theorem's claim for $s \ge 3$ is a claim about the Mellin/Abel-regularised value,
which is what Ramanujan's master theorem computes. This is the same caveat as
Remark~\ref{rem:termwise} one level up: the identity is about the analytic object,
not about any convergent sum or integral of pieces.
\end{remark}

\begin{corollary}[the reading]\label{cor:reading}
Write $\Gamma(\tfrac12)^{2-s}$ as $\Gamma(\tfrac12)^2 \cdot \Gamma(\tfrac12)^{-s}$. The
$2$ is geometry --- the reflection formula at the self-dual exponent forced by the
pair container --- and the $-s$ is combinatorics: \emph{each $S_n$-quotient taken
costs exactly one factor of $\Gamma(\tfrac12)$.} The $\List \to \Bag$ step of \S3 is
$s = 0 \to 1$.
\end{corollary}

This is the kind of statement the programme box of \texttt{q-container.tex} \S``The
programme'' asks for: a coefficient that looks arbitrary at the analytic level
($\pi$ here, $\sqrt{\pi}$ there) forced at the container level, by counting
quotients.

\section{Two obstructions}

Neither of the following is a gap in Theorem~\ref{thm:main}, which is a theorem about
series. Both are gaps in the attempt to make the \emph{whole} statement live at the
container level, and both should be recorded before the reading of
Corollary~\ref{cor:reading} is believed to be more than suggestive.

\subsection{$\Poly$ has no additive inverses, so $-x^2$ is not an object}

$\Poly$, the category of polynomial endofunctors on $\Set$ under $\sqcup$ and
$\otimes$, is a rig category: there is no additive inverse of $y \otimes y$, and
$-x^2$ names no object. Likewise $\Cont_q$. The substitution $F \mapsto F(-x^2)$ is
therefore not container composition $\con$ but composition in the ring of
\emph{virtual} species \cite{Joyal86,Yeh}, that is in the Grothendieck ring $K_0$ of
the species category, where formal differences exist.

\begin{question}\label{q:k0}
Does $\Cont_q$ have a usable $K_0$? Specifically: is the Grothendieck group of
$q$-containers under $+$ (coproduct) a ring under the product of
\texttt{q-container.tex} Def.~\textup{(}product of containers\textup{)}, and does the $q$-series map
$f_C(x) = \sum_s x^{\dim P s}$ extend to it as a ring homomorphism into
$\ZZ[[x]]$?
\end{question}

\begin{warning}\label{warn:k0}
This is the same shape of failure as the divergence section of
\texttt{q-container.tex}: containers resisted Hall convolution because the
construction needed is a convolution on a stack rather than a Day convolution, and
they resist subtraction because $\Poly$ is a rig and not a ring. The two should not
be assumed independent. In particular, do not take the alternating signs of
$F_s(-x^2)$ as evidence of anything container-theoretic until
Question~\ref{q:k0} is settled.
\end{warning}

\subsection{$\int_\RR$ is a second decategorification, not a container operation}

The table of \texttt{q-container.tex} \S``decategorification'' lists one operation
per rung and every one of them is algebraic: coproduct, product, Day convolution,
Hall product, derivative $\partial$ (hyperplanes, decategorifying to the Jackson
$D_q$), shift $\sigma$, $\mathrm{Ind}/\mathrm{Res}$. The species derivative is
pointing --- the one-hole context, the zipper --- and it has a container meaning.
Integration in the sense used here does not: $\int_\RR$ is a measure-theoretic
functional on the middle rung, not the formal antiderivative, and the composite
\[
  \text{container} \;\longrightarrow\; \text{series} \;\longrightarrow\; \text{number}
\]
factors the decategorification into two stages of which only the first is structural.

The nearest precedent is Spivak. In \cite{SpivakEntropy} and
\cite{HosgoodSpivak} a Dirichlet polynomial is read as a set-theoretic bundle, hence
as an empirical distribution, and assigned a Shannon entropy --- a number extracted
from a polynomial functor by a functional that is not one of $\Poly$'s own
operations, and which is nonetheless functorial where it needs to be. That is the
same shape of move as the one made here, and it is the right citation for it.

\begin{conjecture}[\emph{ours}, unverified]\label{conj:compat}
There is a compatibility statement of the following form. Let $\mathcal{B}$ denote
the Borel/Laplace transform of \eqref{eq:borel}, viewed as the de-quotienting map
from symmetric to rigid structures, and let
$\Phi(F) = \int_\RR F(-x^2)\,dx$. Then
$\Phi \circ \mathcal{B} = \Gamma(\tfrac12)\cdot\Phi$ on the relevant
domain; Theorem~\ref{thm:main} is the restriction of this to the family $F_s$.
\end{conjecture}

If Conjecture~\ref{conj:compat} is right it is more than the entropy story has: a
functional that is \emph{not} a container operation nevertheless intertwining with a
map that \emph{is} one, by a single scalar. That scalar is $\Gamma(\tfrac12)$, and
the reason it is $\Gamma$ is that $\Gamma$ is where $n!$ lives.

\section{Next}

\begin{enumerate}[nosep]
  \item \textbf{Settle Question~\ref{q:k0}.} This is the gate. Until $K_0$ of
        $\Cont_q$ exists there is no $q$-analogue of anything in this note, because
        the minus sign is unavailable.
  \item \textbf{The $q$-analogue of Theorem~\ref{thm:main}.} The $q$-side of
        $F_s$ should be $\sum z^n/\norm{\GL_n(\FF_q)}^s$ --- the divisor of the
        $q$-Dirichlet series of \texttt{q-dirichlet.tex} --- and the question is
        what replaces $\Gamma(\tfrac12)$. Jackson's $q$-integral is the obvious
        candidate for $\int$, and $\Gamma_q$ for $\Gamma$; whether the self-dual
        point of Remark~\ref{rem:m2} survives $q$-deformation is the first thing
        to check and is cheap.
  \item \textbf{Prove or refute Conjecture~\ref{conj:compat}} on a domain large
        enough to include a non-hypergeometric example. One such would settle
        whether the statement is structural or an artefact of the family.
  \item \textbf{Add the $(\NN^\delta,+)$ row} of \S2 to
        \texttt{q-dirichlet.tex}, with the observation of
        Proposition~\ref{prop:row}: its three rows all remember $S_n$ and the
        fourth does not, which is the distinction the rest of this note is about.
\end{enumerate}

\begin{thebibliography}{9}

\bibitem{BLL} Fran\c{c}ois Bergeron, Gilbert Labelle and Pierre Leroux,
\emph{Combinatorial Species and Tree-like Structures}, Encyclopedia of Mathematics
and its Applications 67, Cambridge University Press, 1998.

\bibitem{HosgoodSpivak} Tim Hosgood and David I.\ Spivak, \emph{Dirichlet
polynomials and entropy}, arXiv:2107.04832.

\bibitem{Joyal81} Andr\'e Joyal, \emph{Une th\'eorie combinatoire des s\'eries
formelles}, Adv.\ Math.\ \textbf{42} (1981), 1--82.

\bibitem{Joyal86} Andr\'e Joyal, \emph{Foncteurs analytiques et esp\`eces de
structures}, in \emph{Combinatoire \'enum\'erative}, Lecture Notes in Math.\ 1234,
Springer, 1986, 126--159. Virtual species.

\bibitem{Kock} Joachim Kock, \emph{Data types with symmetries and polynomial
functors over groupoids}, Proc.\ MFPS 28, ENTCS 286 (2012), 351--365.

\bibitem{SpivakEntropy} David I.\ Spivak, \emph{Polynomial functors and Shannon
entropy}, arXiv:2201.12878.

\bibitem{Yeh} Yeong-Nan Yeh, \emph{The calculus of virtual species and $K$-species},
in \emph{Combinatoire \'enum\'erative}, Lecture Notes in Math.\ 1234, Springer,
1986, 351--369.

\end{thebibliography}

\end{document}
